\documentclass[a5paper,10pt]{article}

% https://github.com/InDevRus/filippov-solutions

% Нумерация страниц
\pagenumbering{gobble}

% Настройка полей
\usepackage{geometry}
\geometry{left=1cm}
\geometry{right=1cm}
\geometry{top=1cm}
\geometry{bottom=1.5cm}

% Вставка таблицы
\usepackage{array}
\usepackage{tabularx}

% Инструменты выравнивания формул
\usepackage{amsmath}
\usepackage{mathtools}

% Диакритика в математике
\usepackage{amsfonts}

% Вычисления размеров на ходу
\usepackage{calc}

% Удобные записи производных
\usepackage[italic=true]{derivative}

% Настройки сносок
\usepackage{enotez}

% Длинные знаки векторов
\usepackage{anyfontsize}
\usepackage[b]{esvect}

% Вставка картинок
\usepackage{graphicx}

% Вставка красной строки
\usepackage{indentfirst}
\setlength{\parindent}{1em}

% Условия в командах
\usepackage{xifthen}

% Луа-скрипты
\usepackage{luacode}

% Настройка команд отдельно в математическом и текстовом режимах
\usepackage{mathcommand}

% Для повторения LaTeX кода
\usepackage{multido}

% Конфигурация отступов формул
\delimitershortfall-1sp
\usepackage{mleftright}
\mleftright

% Растягивание объектов
\usepackage{relsize}

% Обтекание текстом
\usepackage{wrapfig}
\usepackage{subcaption}
\usepackage{floatflt}

% Поддержка ввода кириллицы
\usepackage[english, russian]{babel}

% Настройка корневой позиции для компилятора
\usepackage{import}
\providecommand*{\ProjectRootPath}{..}

\import{\ProjectRootPath/preambles/macros/}{theorems}

\import{\ProjectRootPath/preambles/fonts}{font_setup}

% Знак градуса
\usepackage{siunitx}

\import{\ProjectRootPath/preambles/macros/}{operators}
\import{\ProjectRootPath/preambles/macros/}{redefinitions}

\import{\ProjectRootPath/preambles/fonts}{letters_setup}
\import{\ProjectRootPath/preambles/fonts/adjustments}{custom_superscripts}

\import{\ProjectRootPath/preambles/custom}{formatting}
\import{\ProjectRootPath/preambles/custom}{frames}
\import{\ProjectRootPath/preambles/custom}{list_setup}

% TODO: перевести команды на современный стиль объявления

\begin{document}

Воспользуемся критерием Михайлова для характеристического уравнения 
$$\pow{\lambda}{4} + 3\pow{\lambda}{3} + 26\pow{\lambda}{2} + 74\lambda + 85 = 0.$$
$p\br{\xi} = \pow{\xi}{2} - 26\xi + 85$,
$q\br{\eta} = -3\eta + 74$.
$\sub{\eta}{1} = \dfrac {74} {3} = 24 \dfrac {2}{3}$.
Критерий $\sub{\xi}{1} < \sub{\eta}{1} < \sub{\xi}{2}$ 
выполнен тогда и только тогда, когда $p\br{\sub{\eta}{1}} < 0$.
При этом квадратичная функция $p\br{\xi}$ имеет точку минимума 
$\sub{\xi}{0} = 13$,
значит правее этой точки $p\br{\xi}$ возрастает и
$p\br{\sub{\eta}{1}} > p\br{24} = 37 > 0$. 
Следовательно, критерий Михайлова не выполняется.

Если $\sub{\lambda}{m} = ai$ -- корень уравнения, 
то тогда $\pow{a}{4} - 3\pow{a}{3} i - 26\pow{a}{2} + 74 ai + 85 = 0$,
откуда
$\tsys{&\ \pow{a}{4} - 26\pow{a}{2} + 85 = 0 
    \\& -3\pow{a}{3} + 74a = 0}$, что равносильно
$\wsys{& \,p\br{\pow{a}{2}} = 0 \\& \,q\br{\pow{a}{2}} = 0}$.
Последнее невозможно, так как уже показано, что эти многочлены общих нулей не имеют.

Значит, чисто мнимых корней нет и обязательно для некоторого корня 
$\Re\br{\lambda} > 0$,
и нулевое решение неустойчиво.
\end{document}
